\begin{tikzpicture}[font=\scriptsize, >=latex]
% O0 vs O2 side by side (real clang riscv64 output counts)
\node[align=left, font=\footnotesize\bfseries] at (4.7,4.6)
  {\ttfamily grade()\rmfamily\ switch: -O0 vs -O2 (real clang 23,
   riscv64 target)};
\draw[thick, fill=graybg] (0.4,1.0) rectangle (4.4,4.2);
\node[font=\footnotesize] at (2.4,3.8) {-O0};
\node[align=left, font=\tiny, text width=3.7cm] at (2.4,2.5)
  {32 instructions:\\frame setup, stores to\\
stack, one branch per case,\\
reload return value, epilogue};
\draw[thick, fill=goodbg] (5.4,1.0) rectangle (9.4,4.2);
\node[font=\footnotesize] at (7.4,3.8) {-O2};
\node[align=left, font=\tiny, text width=3.7cm] at (7.4,2.5)
  {9 instructions, no branches:\\
\texttt{slli, slli, sltiu, add,\\
addiw 60, neg, and, ret} ---\\
the switch became arithmetic};
\node[align=left, text=slate, text width=10cm] at (5.0,0.2)
  {The jump table never appeared: clang proved the four cases were
   linear in n and replaced lookup by computation. Reading
   \ttfamily -O2\rmfamily\ output is how you learn what your C
   really means.};
\end{tikzpicture}
